\section{Planes de ejecución}
La evidencia prevista es \texttt{explain("executionStats")} y
\texttt{EXPLAIN (ANALYZE, BUFFERS)}. Solo se utilizan artefactos existentes;
no se ejecutan consultas para completarlos.
\input{\artifactdir/planes.tex}

\subsection{Escaneos y cardinalidades}
Escaneo secuencial, índice y filtro representan trabajo distinto.
Filas estimadas y reales, loops y cardinalidades intermedias aportan
contexto. La existencia del índice no acredita su elección.

En Q1 desanidar puede ampliar el flujo hacia filtro y agregación.
El total de ítems no determina grupos ni costo del operador.
En Q2 la selectividad condiciona el ordenamiento sin demostrar una
ruta específica.

\subsection{Agregación y buffers}
En PostgreSQL, método de ordenamiento y bloques temporales describen el
plan cuando están presentes. Buffers no son RSS y costos estimados no
son tiempos observados \cite{pgexplain}. Ausencia de campo no significa cero.

En MongoDB, documentos y claves examinados describen acceso. No se
atribuyen etapas ausentes ni se trasladan campos entre versiones.

\subsection{Límites de atribución}
Un EXPLAIN adicional ejecuta su propio trabajo y no completa spill de
otra ventana. El estado CSV no identifica operador ni tamaño temporal.
Un plan Q1 no demuestra Q2 o uso del índice para el cast temporal.
